%%
%% part4-advanced/ch25-custom-commands.tex
%%
%% Chapter 25: Custom Commands — How to define new commands
%% and environments with \newcommand, \NewDocumentCommand,
%% and \NewDocumentEnvironment.
%%

\chapter{دستورات سفارشی}
\label{chap:custom-commands}

\section{چرا این موضوع مهم است؟}
\label{sec:custom-commands-why}

در نوشتن کتاب‌های فنی، بارها پیش می‌آید که
یک دستور یا محیط را چندین بار تکرار می‌کنید.
تعریف دستورات سفارشی، این تکرار را حذف
می‌کند و خوانایی سورس را افزایش می‌دهد.

\lr{MatinBook} خودش از دستورات سفارشی زیادی
استفاده می‌کند (مثل \cmd{cref}، \cmd{pnum}،
\cmd{mnote}). در این فصل، یاد می‌گیرید که
چگونه دستورات و محیط‌های سفارشی خود را
تعریف کنید.

\mnote{تعریف دستور سفارشی، یکی از مهارت‌های کلیدی در \lr{LaTeX} است که به شما امکان می‌دهد سورس کتاب را کوتاه‌تر و خواناتر کنید.}

\section{روش‌های تعریف دستور}
\label{sec:custom-commands-methods}

\lr{LaTeX} سه روش اصلی برای تعریف دستور
ارائه می‌دهد:

\begin{enumerate}
    \item \cmd{newcommand}: روش سنتی.
    \item \cmd{renewcommand}: بازنویسی دستور
    موجود.
    \item \cmd{NewDocumentCommand}: روش مدرن
    \lr{LaTeX3}.
\end{enumerate}

\mnote{\lr{LaTeX3} روش‌های مدرن‌تری برای تعریف دستور ارائه می‌دهد که از تعداد متغیر آرگومان‌ها و آرگومان‌های اختیاری پشتیبانی می‌کنند.}

\section{دستور ساده}
\label{sec:custom-commands-simple}

\subsection{با \cmd{newcommand}}
\label{sec:custom-commands-newcommand}

ساده‌ترین روش تعریف دستور، \cmd{newcommand}
است:

\begin{latin}
\begin{minted}{latex}
\newcommand{\mycommand}{|\pc{متن ثابت}|}
\end{minted}
\end{latin}

این دستور، یک دستور بدون آرگومان تعریف
می‌کند. برای استفاده:

\begin{latin}
\begin{minted}{latex}
\mycommand
\end{minted}
\end{latin}

\mnote{اگر دستور از قبل تعریف شده باشد، \lr{\textbackslash newcommand} خطا می‌دهد. برای اجتناب، از \lr{\textbackslash providecommand} استفاده کنید.}

\subsection{با آرگومان}
\label{sec:custom-commands-args}

برای تعریف دستور با آرگومان:

\begin{latin}
\begin{minted}{latex}
\newcommand{\myemph}[1]{\textbf{#1}}
\end{minted}
\end{latin}

استفاده:

\begin{latin}
\begin{minted}{latex}
\myemph{|\pc{متن پررنگ}|}
\end{minted}
\end{latin}

خروجی: \textbf{متن پررنگ}.

\mnote{تعداد آرگومان‌ها با عدد داخل \lr{[]} تعیین می‌شود. حداکثر ۹ آرگومان امکان‌پذیر است.}

\subsection{با آرگومان اختیاری}
\label{sec:custom-commands-optional}

برای تعریف دستور با آرگومان اختیاری:

\begin{latin}
\begin{minted}{latex}
\newcommand{\mybox}[2][matinblue]{%
    \colorbox{#1}{#2}%
}
\end{minted}
\end{latin}

استفاده:

\begin{latin}
\begin{minted}{latex}
\mybox{|\pc{متن پیش‌فرض}|}%
\mybox[matinred]{|\pc{متن قرمز}|}%
\end{minted}
\end{latin}

\mnote{آرگومان اختیاری با \lr{[...]} و مقدار پیش‌فرض داخل \lr{[]} تعریف می‌شود.}

\section{دستور مدرن (\lr{LaTeX3})}
\label{sec:custom-commands-latex3}

روش مدرن \lr{LaTeX3} از \cmd{NewDocumentCommand}
استفاده می‌کند:

\begin{latin}
\begin{minted}{latex}
\NewDocumentCommand{\myemph}{m}{%
    \textbf{#1}%
}
\end{minted}
\end{latin}

\mnote{روش \lr{LaTeX3} ایمن‌تر است و از انواع آرگومان (اجباری، اختیاری، ستاره‌دار) پشتیبانی می‌کند.}

\subsection{انواع آرگومان}
\label{sec:custom-commands-argtypes}

جدول \cref{tab:arg-types} انواع آرگومان در
\cmd{NewDocumentCommand} را نشان می‌دهد.

\begin{matintable}{انواع آرگومان در \lr{NewDocumentCommand}}{tab:arg-types}
\begin{tblr}{
    width = \textwidth,
    colspec = {c X[l] X[l]},
    hlines, vlines,
    colsep = 8pt,
    rowsep = 4pt,
    row{1} = {font=\bfseries},
}
نوع & توضیح & مثال \\
\lr{m} & آرگومان اجباری & \lr{\{متن\}} \\
\lr{o} & آرگومان اختیاری & \lr{[متن]} \\
\lr{O\{default\}} & اختیاری با پیش‌فرض & \lr{[متن]} \\
\lr{s} & ستاره‌ی اختیاری & \lr{*} \\
\lr{v} & آرگومان verbatim & \lr{|متن|} \\
\end{tblr}
\end{matintable}

\mnote{برای دستورات ساده، \lr{\textbackslash newcommand} کافی است. برای دستورات پیچیده، از \lr{\textbackslash NewDocumentCommand} استفاده کنید.}

\subsection{مثال: دستور با ستاره}
\label{sec:custom-commands-star}

\begin{latin}
\begin{minted}{latex}
\NewDocumentCommand{\myemph}{s m}{%
    \IfBooleanTF{#1}%
        {\textbf{#2}}%
        {\textit{#2}}%
}
\end{minted}
\end{latin}

استفاده:

\begin{latin}
\begin{minted}{latex}
\myemph{|\pc{متن مورب}|}%
\myemph*{|\pc{متن پررنگ}|}%
\end{minted}
\end{latin}

\mnote{این دستور با \lr{\textbackslash myemph} متن را مورب و با \lr{\textbackslash myemph*} متن را پررنگ می‌کند.}

\section{تعریف محیط سفارشی}
\label{sec:custom-commands-env}

\subsection{با \cmd{newenvironment}}
\label{sec:custom-commands-newenv}

برای تعریف محیط سفارشی:

\begin{latin}
\begin{minted}{latex}
\newenvironment{myenv}{%
    \begin{quote}\itshape
}{%
    \end{quote}
}
\end{minted}
\end{latin}

استفاده:

\begin{latin}
\begin{minted}{latex}
\begin{myenv}
|\pc{متن داخل محیط}|%
\end{myenv}
\end{minted}
\end{latin}

\mnote{محیط سفارشی با \lr{\textbackslash newenvironment} تعریف می‌شود و دو بخش دارد: کد شروع و کد پایان.}

\subsection{با \cmd{NewDocumentEnvironment}}
\label{sec:custom-commands-newdocenv}

روش مدرن:

\begin{latin}
\begin{minted}{latex}
\NewDocumentEnvironment{myenv}{m}{%
    \begin{quote}\itshape
    \textbf{#1}\par
}{%
    \end{quote}
}
\end{minted}
\end{latin}

استفاده:

\begin{latin}
\begin{minted}{latex}
\begin{myenv}{|\pc{عنوان}|}
|\pc{متن داخل محیط}|%
\end{myenv}
\end{minted}
\end{latin}

\mnote{روش \lr{LaTeX3} از آرگومان‌های محیط پشتیبانی می‌کند.}

\section{دستورات مفید \lr{MatinBook}}
\label{sec:custom-commands-matinbook}

\lr{MatinBook} چند دستور کمکی برای ساده‌سازی
نوشتن متن ارائه می‌دهد:

\begin{matintable}{دستورات کمکی \lr{MatinBook}}{tab:helper-commands}
\begin{tblr}{
    width = \textwidth,
    colspec = {l X[l] X[l]},
    hlines, vlines,
    colsep = 6pt,
    rowsep = 4pt,
    row{1} = {font=\bfseries},
}
دستور & کاربرد & معادل \\
\lr{\textbackslash cmd\{...\}} & نام دستور & \lr{\textbackslash texttt\{\textbackslash ...\}} \\
\lr{\textbackslash env\{...\}} & نام محیط & \lr{\textbackslash texttt\{...\}} \\
\lr{\textbackslash pkg\{...\}} & نام بسته & \lr{\textbackslash texttt\{...\}} \\
\lr{\textbackslash file\{...\}} & نام فایل & \lr{\textbackslash texttt\{...\}} \\
\lr{\textbackslash opt\{...\}} & آپشن & \lr{\textbackslash texttt\{...\}} \\
\lr{\textbackslash pnum\{...\}} & عدد فارسی & \lr{\textbackslash lr\{...\}} \\
\lr{\textbackslash pc\{...\}} & فارسی در کد & \lr{\textbackslash rl\{...\}} \\
\lr{\textbackslash pcc\{...\}} & کامنت فارسی & — \\
\end{tblr}
\end{matintable}

\mnote{دستورات \lr{\textbackslash cmd}، \lr{\textbackslash env}، \lr{\textbackslash pkg}، \lr{\textbackslash file} و \lr{\textbackslash opt} در \lr{main.tex} مستندات تعریف شده‌اند (نه در کلاس).}

\section{مثال‌های کاربردی}
\label{sec:custom-commands-examples}

\subsection{دستور برای نمایش واژه‌های کلیدی}
\label{sec:custom-commands-keyword}

\begin{latin}
\begin{minted}{latex}
\newcommand{\keyword}[1]{\textbf{\textcolor{matinblue}{#1}}}
\end{minted}
\end{latin}

استفاده:

\begin{latin}
\begin{minted}{latex}
\keyword{|\pc{الگوریتم}|}|\pc{یک مفهوم اساسی است.}|%
\end{minted}
\end{latin}

خروجی: \textbf{\textcolor{matinblue}{الگوریتم}}
یک مفهوم اساسی است.

\subsection{دستور برای نمایش معادل لاتین}
\label{sec:custom-commands-latin-equiv}

\begin{latin}
\begin{minted}{latex}
\newcommand{\latinterm}[2]{#1\lr{(#2)}}
\end{minted}
\end{latin}

استفاده:

\begin{latin}
\begin{minted}{latex}
\latinterm{|\pc{الگوریتم}|}{algorithm}
\end{minted}
\end{latin}

خروجی: الگوریتم\lr{(algorithm)}.

\mnote{این دستور برای نمایش واژه‌ی فارسی به‌همراه معادل لاتین آن در پرانتز مفید است.}

\subsection{محیط برای نکات مهم}
\label{sec:custom-commands-important}

\begin{latin}
\begin{minted}{latex}
\NewDocumentEnvironment{important}{m}{%
    \begin{tcolorbox}[
        colback=matinred!10,
        colframe=matinred,
        title=|\rl{#1}|%
    ]
}{%
    \end{tcolorbox}
}
\end{minted}
\end{latin}

\mnote{برای محیط‌های سفارشی پیچیده، از \lr{tcolorbox} استفاده کنید که امکانات بیشتری دارد.}

\section{خطاهای رایج}
\label{sec:custom-commands-mistakes}

\subsection{تعریف دستور موجود}
\label{sec:custom-commands-mistake-exists}

\textbf{مشکل:} اگر دستوری را که از قبل
وجود دارد با \cmd{newcommand} تعریف کنید،
خطا می‌دهد.

\textbf{راه‌حل:} از \cmd{renewcommand} برای
بازنویسی یا \cmd{providecommand} برای
تعریف امن استفاده کنید.

\subsection{تعداد اشتباه آرگومان‌ها}
\label{sec:custom-commands-mistake-args}

\textbf{مشکل:} اگر تعداد آرگومان‌های تعریف‌شده
با تعداد آرگومان‌های استفاده‌شده یکسان
نباشد، خطا می‌دهد.

\textbf{راه‌حل:} تعداد آرگومان‌ها را در
تعریف و استفاده بررسی کنید.

\subsection{فراموش کردن \cmd{makeatletter}}
\label{sec:custom-commands-mistake-makeatletter}

\textbf{مشکل:} اگر دستوری با \lr{@} در نام
تعریف کنید (مثل \cmd{my@command})،
\lr{LaTeX} خطا می‌دهد.

\textbf{راه‌حل:} از \cmd{makeatletter} و
\cmd{makeatother} استفاده کنید، یا از
\lr{@} در نام دستور اجتناب کنید.

\mnote{در فایل‌های \lr{.sty}، استفاده از \lr{@} در نام دستورات داخلی مرسوم است، ولی در سند اصلی باید از \lr{\textbackslash makeatletter} استفاده کنید.}

\section{تمرین‌ها}
\label{sec:custom-commands-exercises}

\begin{exercise}
    یک دستور ساده تعریف کنید که متن را
    پررنگ کند.
    \label{exc:custom-commands-bold}
\end{exercise}

\begin{exercise}
    یک دستور با دو آرگومان تعریف کنید
    که یک واژه و معادل لاتین آن را نمایش
    دهد.
    \label{exc:custom-commands-two-args}
\end{exercise}

\begin{exercise}
    یک دستور با آرگومان اختیاری تعریف کنید
    که رنگ متن را تغییر دهد.
    \label{exc:custom-commands-optional}
\end{exercise}

\begin{exercise}
    یک محیط سفارشی تعریف کنید که متن را
    در یک جعبه‌ی رنگی قرار دهد.
    \label{exc:custom-commands-env}
\end{exercise}

\section{خلاصه}
\label{sec:custom-commands-summary}

در این فصل، با دستورات سفارشی در \lr{MatinBook}
آشنا شدیم:

\begin{itemize}
    \item سه روش تعریف دستور:
    \cmd{newcommand}، \cmd{renewcommand}،
    \cmd{NewDocumentCommand}.

    \item دستور ساده بدون آرگومان:
    \cmd{newcommand\{\textbackslash mycmd\}\{...\}}.

    \item دستور با آرگومان:
    \cmd{newcommand\{\textbackslash mycmd\}[1]\{...\}}.

    \item دستور با آرگومان اختیاری:
    \cmd{newcommand\{\textbackslash mycmd\}[2][default]\{...\}}.

    \item روش مدرن \lr{LaTeX3} از انواع آرگومان
    (اجباری، اختیاری، ستاره‌دار) پشتیبانی
    می‌کند.

    \item محیط سفارشی با \cmd{newenvironment}
    یا \cmd{NewDocumentEnvironment}.

    \item دستورات کمکی \lr{MatinBook}:
    \cmd{cmd}، \cmd{env}، \cmd{pkg}،
    \cmd{file}، \cmd{opt}، \cmd{pnum}،
    \cmd{pc}، \cmd{pcc}.

    \item خطاهای رایج شامل تعریف دستور موجود،
    تعداد اشتباه آرگومان‌ها، و فراموش کردن
    \cmd{makeatletter} است.
\end{itemize}

در فصل بعدی (\cref{chap:large-projects})،
با مدیریت پروژه‌های بزرگ آشنا می‌شوید و
یاد می‌گیرید که چگونه کتاب‌های چندفایلی
را سازمان‌دهی کنید.